\section{案例拆解：Web、代码与研究三类任务}

\subsection{案例一：Web Agent 的长链交互}
在 Web 场景中，Agent 往往需要跨页面导航、解析半结构化信息并完成多步决策。课程观点是，这类任务的失败通常不是 ``模型不懂页面''，而是 ``状态同步与计划更新不及时''。典型问题包括：页面 DOM 变化导致定位失效、登录态中断、同义目标词触发错误点击、以及工具返回值与视觉状态不一致。

\begin{importantbox}{Web Agent 的核心难点是“弱结构环境中的强闭环控制”}
Agent 必须在不稳定的界面状态下持续完成 Observe-Reason-Act 闭环。若缺乏显式状态机和步骤级验证器，错误会跨轮扩散，最终表现为随机性失败。
\end{importantbox}

一个可行策略是把页面操作划分为 ``导航动作''、``提取动作''、``确认动作'' 三类，并为每类动作配置不同验证逻辑。比如导航动作关注 URL/页面标题变化，提取动作关注字段完整性，确认动作关注业务约束是否满足。这样做可以降低 ``一步错全盘错'' 的概率。

\begin{knowledgebox}{Web Agent 评测建议}
\begin{itemize}
    \item 记录动作级成功率，而非只记录终局成功率；
    \item 统计重试次数与回退次数，评估计划稳定性；
    \item 按页面复杂度分桶，避免平均值掩盖难例。
\end{itemize}
\end{knowledgebox}

\subsection{案例二：Code Agent 的仓库级任务}
代码任务看似结构化，但仓库级目标通常涉及跨文件依赖、隐式规范和历史约束。课程强调，Code Agent 的提升路径不只是 ``更强代码生成''，而是 ``更强的仓库状态理解 + 更强的验证反馈回路''。

\begin{warningbox}{只做单轮补丁会放大技术债}
如果 Agent 只追求当前测试通过，而忽略架构一致性、命名规则和接口契约，短期成功会变成长期维护成本。Code Agent 需要把 lint、type check、unit test、integration test 联合作为验证器。
\end{warningbox}

在实现层面，可将任务分为 ``理解''、``修改''、``验证''、``复盘'' 四阶段，并把每阶段输出写入 episodic memory。这样做的价值是：当任务失败时，系统可以定位到底是需求理解偏差、实现错误还是验证覆盖不足，从而触发针对性重规划。

\begin{knowledgebox}{Code Agent 的 memory 建议}
\begin{itemize}
    \item semantic memory 存储项目规范、公共接口和历史决策；
    \item episodic memory 存储本次任务的关键尝试与失败原因；
    \item procedural memory 存储复用脚本（如回归测试工作流）。
\end{itemize}
\end{knowledgebox}

\subsection{案例三：Research Agent 的证据闭环}
研究任务常被误判为 ``纯 reasoning 问题''。Yu Su 的框架下，Research Agent 同样需要强 memory 与 planning：前者用于证据沉淀与冲突管理，后者用于研究路线切换和实验预算分配。

\begin{importantbox}{Research Agent 的关键不是“写得像论文”，而是“证据可追溯”}
系统应为每条结论绑定来源、置信度和时间戳。否则，当检索语料更新或工具返回变化时，旧结论会在记忆层中持续污染后续推理。
\end{importantbox}

一个实用做法是采用 ``结论-证据-反证'' 三元结构写入 memory：Agent 每次提出结论时必须附上支持证据与潜在反例。这能显著提升反思（reflection）阶段的有效性，降低 ``解释型幻觉''。

\begin{table}[H]
\centering
\begin{tabular}{p{2.7cm}p{3.9cm}p{4.4cm}}
\toprule
\textbf{任务类型} & \textbf{优先能力} & \textbf{首要风险} \\
\midrule
Web Agent & 状态同步 + 动态重规划 & 页面漂移导致动作失效 \\
Code Agent & 验证闭环 + 仓库记忆 & 局部最优补丁侵蚀全局质量 \\
Research Agent & 证据管理 + 反思机制 & 无来源结论的记忆污染 \\
\bottomrule
\end{tabular}
\caption{三类代表任务的能力优先级差异}
\end{table}

\subsection{本章小结}
三类任务的共同点是都需要系统闭环；差异在于验证器和记忆策略的设计重点不同。Agent 研发应按任务属性配置能力，而非追求单一通用模板。

